\chapter{Conclusion and Future Work}\label{chap:conclusion}
We have presented \Rtrace{}, an operational semantics for helping React developers understand the behavior of Hooks.
Render-related UI bugs often puzzle developers due to the peculiar semantics of Hooks.
We have captured the essence of Hooks in \Rtrace{}, which clearly explains subtle pitfalls that occur when multiple Hooks interact with each other.
Moreover, developers can use our \Rtrace-interpreter and visualizer to inspect how their programs render step-by-step, making the opaque render process transparent.

As a foundation for future work, \Rtrace{} opens several research directions.
First, our semantics enables semantic-based static analyzers and type systems that go beyond the syntactic or runtime checks of current tools, catching subtle bugs at compile time.
Second, as React evolves with features like React compiler and server components, \Rtrace{} provides a formal basis for verifying their correctness.
Third, \Rtrace{} can be refined for pedagogical use in building a correct mental model of Hooks.
Finally, integrating \Rtrace{} with web semantics can enable end-to-end verification of web programs, from component behavior to final rendering.
